\documentclass[11pt,letterpaper]{article}
\usepackage[T1]{fontenc}
\usepackage{lmodern}
\usepackage[margin=1in,headheight=15pt]{geometry}
\usepackage{microtype}
\usepackage{amsmath,amssymb,amsthm,mathtools}
\usepackage{booktabs,array,longtable}
\usepackage[dvipsnames]{xcolor}
\usepackage{enumitem}
\usepackage{fancyhdr}
\usepackage{titlesec}
\usepackage{listings}
\usepackage{hyperref}
\definecolor{accent}{HTML}{45634B}
\definecolor{ink}{HTML}{24302A}
\definecolor{soft}{HTML}{F2F4ED}
\hypersetup{colorlinks=true,linkcolor=accent,citecolor=accent,urlcolor=accent,
  pdftitle={Product formulas for ballot-polynomial Hankel determinants},
  pdfauthor={Research report prepared for Vladimir Reshetnikov},
  pdfsubject={Proofs of Cigler's Conjectures 13--15 and strict log-concavity}}
\titleformat{\section}{\large\bfseries\color{accent}}{\thesection}{0.7em}{}
\titleformat{\subsection}{\normalsize\bfseries\color{ink}}{\thesubsection}{0.7em}{}
\pagestyle{fancy}
\fancyhf{}
\fancyhead[L]{\small Ballot-polynomial Hankel determinants}
\fancyhead[R]{\small Research note}
\fancyfoot[C]{\thepage}
\renewcommand{\headrulewidth}{0.3pt}
\setlength{\parindent}{1.2em}
\setlength{\parskip}{3pt}
\setlength{\emergencystretch}{2em}
\numberwithin{equation}{section}
\newtheorem{theorem}{Theorem}[section]
\newtheorem{lemma}[theorem]{Lemma}
\newtheorem{proposition}[theorem]{Proposition}
\newtheorem{corollary}[theorem]{Corollary}
\theoremstyle{definition}
\newtheorem{definition}[theorem]{Definition}
\newtheorem{example}[theorem]{Example}
\newtheorem{remark}[theorem]{Remark}
\newcommand{\Z}{\mathbb Z}
\newcommand{\Q}{\mathbb Q}
\newcommand{\N}{\mathbb N}
\newcommand{\eps}{\varepsilon}
\newcommand{\ct}[2]{[#1]#2}
\newcommand{\fall}[2]{#1^{\underline{#2}}}
\newcommand{\rising}[2]{(#1)_{#2}}
\DeclareMathOperator{\diag}{diag}
\DeclareMathOperator{\ord}{ord}
\lstset{basicstyle=\small\ttfamily,breaklines=true,columns=fullflexible,
  frame=single,rulecolor=\color{black!15},backgroundcolor=\color{soft},
  keepspaces=true,showstringspaces=false}

\begin{document}
\thispagestyle{empty}
\begin{center}
{\color{accent}\small DISCRETE MATHEMATICS \quad / \quad EXACT DETERMINANT EVALUATION}\par
\vspace{0.6em}
{\LARGE\bfseries Product formulas for\par
ballot-polynomial Hankel determinants}\par
\vspace{0.6em}
{\large Proofs of Cigler's Conjectures 13--15,\par
with strict log-concavity and coefficient asymptotics}\par
\vspace{1em}
{\normalsize A research report prepared for Vladimir Reshetnikov}\par
{\small 19 September 2026}
\end{center}
\vspace{0.5em}
\begin{abstract}
We study the shifted Hankel determinants of the ballot polynomials
\(a_m(t)=\sum_{h\le m/2}(\binom mh-\binom m{h-1})t^h\).
An explicit evaluation expresses every normalized determinant as a sum of
only \(\lfloor k/2\rfloor+1\) monomials, whose coefficients are products
of two elementary Vandermonde-type factors. This establishes the sign,
support, and positivity asserted in Conjecture 13 of Cigler's 2021 paper.
The same products are polynomial in the matrix size on each parity class.
Their degree and a negative-index reflection establish the rational
generating functions, exact numerator degrees, and reciprocity identities
of Conjectures 14 and 15. Moreover, all the residual coefficient sequences
are strictly log-concave, and after normalization they converge to a
binomial distribution for fixed shift. The proof is self-contained: a
block determinant reduces an \(n\)-by-\(n\) Hankel determinant to a
\(k\)-by-\(k\) coefficient determinant; a two-term orthogonal-polynomial
connection then collapses its expansion to parity-separated Vandermondes.
Exact verification code, symbolic numerator data, and coefficient tables
accompany the article.
\end{abstract}

\section{The selected problem and the outcome}\label{sec:problem}
The problem is a small, sharply delimited part of Johann Cigler's paper
\emph{Hankel determinants of middle binomial coefficients and conjectures
for some polynomial extensions and modifications}
\cite{Cigler2021}. Section 5, pages 19--21 of version 3, poses three
conjectures about the same family of determinants. Their combination is
particularly suitable for a direct attack: the moments have a very simple
orthogonal-polynomial recurrence, and the predicted support of the
normalized determinant is strikingly small.

For integers \(m\ge0\), put
\begin{equation}\label{eq:ballot}
 a_m(t)=\sum_{h=0}^{\lfloor m/2\rfloor}
 \left(\binom mh-\binom m{h-1}\right)t^h,
 \qquad \binom m{-1}=0.
\end{equation}
The first polynomials are
\[
 1,\quad 1,\quad 1+t,\quad 1+2t,\quad
 1+3t+2t^2,\quad 1+4t+5t^2.
\]
At \(t=1\), the sum telescopes to
\(a_m(1)=\binom m{\lfloor m/2\rfloor}\), the middle binomial
coefficients, OEIS A001405 \cite{OEIS1405}.
Define
\begin{equation}\label{eq:hankel}
 \Delta_k(n,t)=\det\bigl(a_{k+i+j}(t)\bigr)_{0\le i,j<n},
 \qquad
 \delta_k(n,t)=\frac{\Delta_k(n,t)}{t^{\binom n2}}.
\end{equation}
The empty determinant is 1. We first interpret the quotient over
\(\Q(t)\); the proof below shows that it is a polynomial, so its value
at \(t=0\) is obtained by polynomial specialization, not division by zero.

Write \(F_k(x,t)=\sum_{n\ge0}\delta_k(n,t)x^n\).
Here is the exact scope of the three published conjectures, in our indexing.
For \(k\ge1\), set
\begin{equation}\label{eq:e-d}
 e_k(n)=\left\lfloor\frac{k(n-1)+1}{2}\right\rfloor,
 \qquad r=\lfloor k/2\rfloor,
 \qquad d_k=\lfloor k^2/4\rfloor.
\end{equation}
Conjecture 13 predicts
\begin{equation}\label{eq:c13}
 \delta_k(n,t)=(-1)^{k\binom n2}t^{e_k(n)}\rho_k(n,t),
 \qquad \deg_t\rho_k(n,t)=r,
\end{equation}
with positive coefficients. For \(n\ge1\), we prove positivity of
\emph{every} coefficient of \(t^0,\ldots,t^r\).
At the empty-matrix boundary, \(\rho_k(0,t)=t^r\); it is best to state
that boundary separately rather than include it in a strict-positivity
assertion.

Conjecture 14 predicts
\begin{align}
 F_{2r}(x,t)&=
 \frac{A_{2r}(x,t)}{(1-t^r x)(1-t^{2r}x^2)^{r^2}},
 &\deg_x A_{2r}&=2r(r-1)+1, \quad r\ge1,\label{eq:c14even}\\
 F_{2r+1}(x,t)&=
 \frac{A_{2r+1}(x,t)}{(1+t^{2r+1}x^2)^{r^2+r+1}},
 &\deg_x A_{2r+1}&=2r^2+1, \quad r\ge0.\label{eq:c14odd}
\end{align}
Conjecture 15 predicts the two reciprocity identities
\begin{align}
 x^{2r(r-1)+1}t^{2r^2(r-1)}
 A_{2r}(x^{-1},t^{-1})&=A_{2r}(x,t),\label{eq:c15even}\\
 (-1)^r x^{2r^2+1}t^{(2r+1)r^2}
 A_{2r+1}(x^{-1},t^{-1})&=A_{2r+1}(x,t).\label{eq:c15odd}
\end{align}
We prove all these formulas. The prescribed denominators are not asserted
to remain reduced after every specialization of \(t\).

\paragraph{Status and attribution.}
The cited source is arXiv:2111.14492v3, revised 30 December 2021, where
these statements are explicitly labeled conjectures. A targeted literature
search conducted for this report, through 19 September 2026, did not locate
a subsequent proof of these particular Section 5 statements. That search
is not a certification of priority. The mathematical claim here is that
the detailed argument below proves the displayed statements; the report
has not undergone external peer review or proof-assistant verification.
Orthogonal-polynomial reductions of Hankel determinants are established
techniques, not a claimed invention of this report; related general
results appear in \cite{CK2020,K2021}. The useful specialization here is
the collapse to one deleted column and the resulting explicit coefficient
products.

\paragraph{Road map.}
Section~\ref{sec:product} gives the product evaluation. Sections
\ref{sec:moments}--\ref{sec:proof} prove it from first principles.
Section~\ref{sec:logconcavity} strengthens positivity to strict
log-concavity. Sections~\ref{sec:poly}--\ref{sec:reciprocity} derive the
polynomial dependence, asymptotics, generating functions, and reciprocity.
The final sections discuss computation, examples, and reproducibility.

\section{An explicit coefficient product}\label{sec:product}
It is convenient to keep a second weight:
\begin{equation}\label{eq:two-weights}
 a_m(u,t)=\sum_{h=0}^{\lfloor m/2\rfloor}
 \left(\binom mh-\binom m{h-1}\right)u^{m-2h}t^h.
\end{equation}
Define \(\Delta_k(n;u,t)\) and \(\delta_k(n;u,t)\) as in
\eqref{eq:hankel}, with these moments. Setting \(u=1\) recovers the
original problem.

\begin{definition}\label{def:coeff}
For \(k\ge1\) and an integer \(n\), let
\begin{equation}\label{eq:admissible}
 \mathcal A_{k,n}=
 \begin{cases}
  \{0,2,\ldots,k\},&k\text{ even},\\
  \{a:0\le a\le k,\ a\equiv n\pmod2\},&k\text{ odd}.
 \end{cases}
\end{equation}
For \(a\in\mathcal A_{k,n}\), delete one integer from a consecutive
interval:
\[
 S_{k,n,a}=\{n-1,n,\ldots,n+k-1\}\setminus\{n+a-1\}.
\]
From its even and odd elements form the increasing lists
\[
 E=\bigl(d/2:d\in S_{k,n,a},\ d\text{ even}\bigr),
 \qquad
 O=\bigl((d-1)/2:d\in S_{k,n,a},\ d\text{ odd}\bigr).
\]
Their lengths are \(e=\lceil k/2\rceil\) and \(o=\lfloor k/2\rfloor\).
Define
\begin{align}
 V_E(E)&=
 \frac{\displaystyle\prod_{p<q}(E_q-E_p)(E_q+E_p+1)}
 {\displaystyle\prod_{i=0}^{e-1}(2i)!},\label{eq:VE}\\
 V_O(O)&=
 \frac{\displaystyle\prod_{p=0}^{o-1}(O_p+1)
       \prod_{p<q}(O_q-O_p)(O_q+O_p+2)}
 {\displaystyle\prod_{i=0}^{o-1}(2i+1)!},\label{eq:VO}\\
 C_{k,a}(n)&=V_E(E)V_O(O).\label{eq:C}
\end{align}
All products over empty sets are 1.
\end{definition}

\begin{theorem}[Product evaluation]\label{thm:main}
For \(k,n\ge1\), every \(C_{k,a}(n)\) is a positive integer, and
\begin{equation}\label{eq:main}
 \boxed{\displaystyle
 \delta_k(n;u,t)=(-1)^{k\binom n2}
  \sum_{a\in\mathcal A_{k,n}}
   C_{k,a}(n)\,u^a\,t^{(kn-a)/2}.}
\end{equation}
In particular, this normalized determinant contains exactly
\(\lfloor k/2\rfloor+1\) nonzero monomials before specializing \(u\).
\end{theorem}

To read off \(\rho_k\), put
\[
 \eps=\begin{cases}1,&k\text{ and }n\text{ both odd},\\0,&\text{otherwise}.
 \end{cases}
 \qquad
 b_j=C_{k,\eps+2j}(n)\quad(0\le j\le r).
\]
Then
\begin{equation}\label{eq:rho-b}
 \rho_k(n,t)=\sum_{j=0}^r b_j t^{r-j}.
\end{equation}
Indeed, the least exponent in \eqref{eq:main} is
\((kn-\eps-2r)/2=e_k(n)\), and the largest is
\((kn-\eps)/2=\lfloor kn/2\rfloor\).
Consequently Theorem~\ref{thm:main} proves Conjecture 13, including both
ends of the support; it proves more than a bound on the degree.

\begin{example}\label{ex:k4n2}
Take \(k=4\), \(n=2\). The degree interval is \(1,2,3,4,5\), and
\(a=0,2,4\) deletes respectively \(1,3,5\). In all three cases
\(E=(1,2)\), so \(V_E=2\). The corresponding odd lists are
\((1,2),(0,2),(0,1)\), giving \(V_O=5,4,1\). Thus
\[
 (b_0,b_1,b_2)=(10,8,2),\qquad
 \delta_4(2;u,t)=10t^4+8u^2t^3+2u^4t^2.
\]
Since \(e_4(2)=2\), the residual polynomial is
\(\rho_4(2,t)=2+8t+10t^2\).
\end{example}

\section{The moments and their orthogonal polynomials}\label{sec:moments}
\subsection{A weighted path model and the moment generating function}
Consider nonnegative lattice paths that start and end at height 0.
An up-step raises the height by 1 and has weight 1; a down-step lowers it
by 1 and has weight \(t\). Horizontal steps are permitted only at height
0 and have weight \(u\). Length is marked by \(w\). The unweighted
interpretation is the dispersed-Dyck-path model associated with the middle
binomial coefficients \cite{Cigler2021,OEIS1405}.

Let \(C(v)\) be the Catalan generating function, uniquely determined by
\(C(v)=1+vC(v)^2\), \(C(0)=1\). A path is a sequence of either a
horizontal step or a positive excursion. A positive excursion contributes
\(tw^2C(tw^2)\). Therefore its generating function is
\begin{align}
 M(w;u,t)&=\frac{1}{1-uw-tw^2C(tw^2)}\notag\\
 &=\frac{C(tw^2)}{1-uwC(tw^2)}
 =\frac{2}{1-2uw+\sqrt{1-4tw^2}}.\label{eq:M}
\end{align}
For completeness, the identity \(1-vC(v)=1/C(v)\) gives the second
expression directly from the first. Expanding it in powers of \(u\)
gives
\[
 M(w;u,t)=\sum_{s\ge0}u^sw^s C(tw^2)^{s+1}.
\]
Lagrange inversion applied to \(C-1=vC^2\) gives
\[
 [v^h]C(v)^{s+1}
 =\frac{s+1}{2h+s+1}\binom{2h+s+1}{h}.
\]
With \(m=s+2h\), this is
\(\binom mh-\binom m{h-1}\). Hence
\([w^m]M(w;u,t)=a_m(u,t)\), proving that the path model has exactly
our moments. This also explains the homogeneous weight relation
\(u\mapsto\lambda u\), \(t\mapsto\lambda^2t\).

\subsection{A tridiagonal operator}
Work initially over the field \(\Q(u,t)\). On a vector space with basis
\(e_0,e_1,\ldots\), define
\[
 J e_0=e_1+u e_0,\qquad
 J e_h=e_{h+1}+t e_{h-1}\quad(h\ge1).
\]
For a polynomial \(f\), let
\(L(f)=[e_0]f(J)e_0\). The path interpretation gives
\(L(z^m)=a_m(u,t)\).

Define monic polynomials
\begin{equation}\label{eq:U}
 U_{-1}(z;t)=0,\qquad U_0(z;t)=1,\qquad
 U_m(z;t)=zU_{m-1}(z;t)-tU_{m-2}(z;t)\quad(m\ge1),
\end{equation}
and put
\begin{equation}\label{eq:P}
 P_m(z)=U_m(z;t)-uU_{m-1}(z;t)\quad(m\ge0).
\end{equation}
Thus \(P_0=1\), \(P_1=z-u\), and
\(P_m=zP_{m-1}-tP_{m-2}\) for \(m\ge2\).
These are rescaled Chebyshev polynomials with a one-term perturbation of
the initial condition; no special-function facts are needed below.

An induction gives \(P_m(J)e_0=e_m\): the cases \(m=0,1\) follow
from the definition of \(J\), and the recurrence advances the identity.
Consequently
\begin{equation}\label{eq:orthogonal}
 L(z^iP_m)=0\quad(i<m),\qquad L(z^mP_m)=t^m.
\end{equation}
The second equality follows because the only length-\(m\) path from
height \(m\) to height 0 consists of \(m\) down-steps.
Since the \(P_m\) are monic, \eqref{eq:orthogonal} also gives
\(L(P_iP_j)=0\) for \(i\ne j\) and \(L(P_j^2)=t^j\).
Changing from monomials to the monic \(P_j\)'s in a Gram determinant
therefore proves
\begin{equation}\label{eq:unshifted}
 \Delta_0(n;u,t)=\prod_{j=0}^{n-1}t^j=t^{\binom n2}.
\end{equation}

\subsection{A self-contained small-determinant reduction}
\begin{lemma}\label{lem:reduction}
For \(k,n\ge0\),
\begin{equation}\label{eq:small}
 \delta_k(n;u,t)=(-1)^{nk}
 \det\bigl([z^i]P_{n+j}(z)\bigr)_{0\le i,j<k}.
\end{equation}
\end{lemma}
\begin{proof}
Form an \((n+k)\)-by-\((n+k)\) matrix whose columns are indexed by
\(P_0,\ldots,P_{n+k-1}\). Its first \(n\) rows are the functionals
\(f\mapsto L(z^if)\), \(0\le i<n\), and its last \(k\) rows are
\(f\mapsto[z^i]f\), \(0\le i<k\).
By \eqref{eq:orthogonal}, its upper right block is zero and the upper
left block is triangular with diagonal \(1,t,\ldots,t^{n-1}\).
Thus its determinant is
\[
 t^{\binom n2}\det\bigl([z^i]P_{n+j}\bigr)_{0\le i,j<k}.
\]
Replace the monic-polynomial columns by the monomial columns
\(1,z,\ldots,z^{n+k-1}\). The change-of-basis determinant is 1.
The bottom \(k\) rows are now an identity block in the first \(k\)
columns and zero elsewhere. Expanding along these rows produces the sign
\((-1)^{nk}\). The remaining matrix is
\((L(z^{k+i+j}))_{0\le i,j<n}\), whose determinant is
\(\Delta_k(n;u,t)\). Equating the two expressions and using
\eqref{eq:unshifted} proves the assertion.
\end{proof}
This proof isolates the familiar orthogonal-polynomial reduction rather
than treating it as a black box. The remaining work is a special
evaluation of its right-hand side.

\section{Why only a few minors survive}\label{sec:proof}
\subsection{The two-term connection collapses the expansion}
Expand the determinant in \eqref{eq:small} using
\(P_{n+j}=U_{n+j}-uU_{n+j-1}\). A priori there are \(2^k\)
choices. Call the first choice in a column \emph{unshifted} and the
second \emph{shifted}. If an unshifted column is followed immediately
by a shifted column, both selected polynomials are \(U_{n+j}\), so that
term vanishes. Every binary choice pattern other than a shifted prefix
followed by an unshifted suffix has such an adjacent pair. Therefore
only \(k+1\) terms can survive.

If the shifted prefix has length \(a\), its scalar coefficient is
\((-u)^a\), and its increasing list of polynomial indices is precisely
\(S_{k,n,a}\): it is the interval \(n-1,\ldots,n+k-1\) with
\(n+a-1\) omitted. Each \(U_d\) has the parity of \(d\). There
are \(e=\lceil k/2\rceil\) even coefficient rows and
\(o=\lfloor k/2\rfloor\) odd rows in \eqref{eq:small}. A term vanishes
unless the selected column parities have exactly those multiplicities.
Counting the two parities in the consecutive interval gives exactly
\(a\in\mathcal A_{k,n}\), leaving \(r+1\) terms.

\subsection{Vandermonde evaluation of the two blocks}
The recurrence \eqref{eq:U} implies
\[
 U_d(z;t)=\sum_{h=0}^{\lfloor d/2\rfloor}
  (-t)^h\binom{d-h}{h}z^{d-2h}\qquad(d\ge0).
\]
In particular,
\begin{align}
 [z^{2i}]U_{2m}(z;t)&=(-t)^{m-i}\binom{m+i}{2i},\label{eq:evenentry}\\
 [z^{2i+1}]U_{2m+1}(z;t)&=(-t)^{m-i}\binom{m+i+1}{2i+1}.\label{eq:oddentry}
\end{align}
The binomial factors have useful polynomial forms:
\begin{align}
 \binom{m+i}{2i}
 &=\frac{\displaystyle\prod_{h=0}^{i-1}
       \bigl(m(m+1)-h(h+1)\bigr)}{(2i)!},\label{eq:quadE}\\
 \binom{m+i+1}{2i+1}
 &=\frac{m+1}{(2i+1)!}
    \prod_{h=1}^{i}\bigl((m+1)^2-h^2\bigr).\label{eq:quadO}
\end{align}
For the even block, the row indexed by \(i\), apart from its factorial
factor, is a monic degree-\(i\) polynomial in \(m(m+1)\).
Its determinant is therefore the Vandermonde determinant in those
quadratic arguments. Since
\[
 E_q(E_q+1)-E_p(E_p+1)
 =(E_q-E_p)(E_q+E_p+1),
\]
it evaluates to \(V_E(E)\).
For the odd block, first factor \(m+1\) from each column and apply the
same argument to the squared arguments \((m+1)^2\). This gives
\(V_O(O)\).

This argument also proves integrality without an appeal to divisibility
properties of the product expressions: \(V_E\) and \(V_O\) are
respectively the determinants of the integer matrices
\[
 \left(\binom{E_j+i}{2i}\right)_{0\le i,j<e},\qquad
 \left(\binom{O_j+i+1}{2i+1}\right)_{0\le i,j<o}.
\]
When \(n\ge1\), all nodes are nonnegative, every node difference is
positive, and every sum factor appearing in these products is positive.
The factors \(O_j+1\) are positive as well. Thus both determinants
are positive integers.

Factoring the powers of \((-t)\) in
\eqref{eq:evenentry}--\eqref{eq:oddentry} contributes the exponent
\[
 W=\sum E_j+\sum O_j-\binom e2-\binom o2.
\]
The selected degrees have sum
\(kn+k(k-1)/2-a\). Since their sum also equals
\(2\sum E_j+2\sum O_j+o\), substitution gives
\begin{equation}\label{eq:W}
 W=(kn-a)/2.
\end{equation}

\subsection{The common sign}
The positivity calculation must not hide the determinant signs. Let
\(R\) be the number of odd-before-even pairs in the row order
\(0,1,\ldots,k-1\), and let \(I(S)\) be the corresponding number
in the increasing selected degree list. Permuting both lists into even
and odd blocks has sign \((-1)^{R+I(S)}\). Including
\eqref{eq:small}, the prefix coefficient \((-u)^a\), and
\eqref{eq:W}, the overall sign exponent is
\begin{equation}\label{eq:signexp}
 nk+a+(kn-a)/2+R+I(S).
\end{equation}
When \(a\) is increased by 2 within the admissible set, the selected
parity word changes by interchanging one adjacent even--odd pair.
Thus \(I(S)\) changes parity. At the same time \((kn-a)/2\)
decreases by 1 and \(a\) increases by 2. The total sign is unchanged.
It suffices to check the smallest admissible \(a\).

For even \(n\), that index is 0 and both parity words begin with an
even element, so \(I(S)=R\). Formula \eqref{eq:signexp} then has
parity \(kn/2\), which agrees with \(k\binom n2\).
For odd \(n\), the two cases are
\[
\begin{array}{c|c|c|c}
 k & a_{\min} & R & I(S)\\ \hline
 2r & 0 & r(r-1)/2 & r(r+1)/2\\
 2r+1 & 1 & r(r+1)/2 & r(r-1)/2.
\end{array}
\]
For \(k=2r\), \(R+I(S)=r^2\) and \((kn-a)/2=rn\), so
the resulting sign is positive, as required. For \(k=2r+1\),
\eqref{eq:signexp} has parity \((n-1)/2\), which is the parity of
\(k\binom n2\). We have therefore proved that every surviving term
has the sign \((-1)^{k\binom n2}\).
Together with the product evaluation and \eqref{eq:W}, this completes
the proof of Theorem~\ref{thm:main}.

\section{Strict log-concavity: a stronger conclusion}\label{sec:logconcavity}
The product evaluation gives more than positivity. The dependence on the
deleted node can be made almost one-dimensional.

\begin{theorem}[Adjacent ratios]\label{thm:ratios}
Let \(k,n\ge1\), \(r=\lfloor k/2\rfloor\), and
\(b_j=C_{k,\eps+2j}(n)\). If it is not the case that \(k\) is even
and \(n\) is odd, put \(L=n+(n\bmod2)\). Then
\begin{equation}\label{eq:ratioRegular}
 \frac{b_{j+1}}{b_j}
 =\frac{(r-j)(L+j)}{(j+1)(L+r+j+1)}\quad(0\le j<r).
\end{equation}
In the remaining case, \(k\) even and \(n\) odd,
\begin{equation}\label{eq:ratioExceptional}
 \frac{b_{j+1}}{b_j}
 =\frac{(r-j)(n+j)(n+2j+2)}
 {(j+1)(n+r+j+1)(n+2j)}.
\end{equation}
\end{theorem}
\begin{proof}
All dependence on \(j\) occurs in the parity block from which one
node is deleted. In the first set of cases that block is the odd block.
Its full node list, before deletion, can be written
\(m+1=L/2+\ell\), \(0\le\ell\le r\).
Deleting \(\ell=j\) removes one column factor and every
Vandermonde factor involving that column. Ignoring a positive factor
independent of \(j\), the remaining determinant equals
\begin{equation}\label{eq:bRegular}
 b_j=\frac{K}{j!(r-j)!\,\rising{L+j}{r+1}},
 \qquad K>0.
\end{equation}
To see the cancellation explicitly, the product of the removed node
sums is
\(\rising{L+j}{r+1}/(L+2j)\), whereas the removed odd column factor
is \(L/2+j\). Their quotient supplies the constant factor 2, which
is absorbed in \(K\). Here \((z)_s=z(z+1)\cdots(z+s-1)\).

For even \(k\) and odd \(n\), the deleted block is instead the even
block, with full nodes \(m=(n-1)/2+\ell\). There is no odd column
factor to cancel the omitted diagonal sum. The corresponding expression is
\begin{equation}\label{eq:bExceptional}
 b_j=\frac{K(n+2j)}{j!(r-j)!\,\rising{n+j}{r+1}},
 \qquad K>0.
\end{equation}
Taking adjacent quotients in \eqref{eq:bRegular} and
\eqref{eq:bExceptional} gives the two stated formulas.
\end{proof}

\begin{corollary}[Strict log-concavity]\label{cor:logconcavity}
For every \(k,n\ge1\), the positive coefficient sequence of
\(\rho_k(n,t)\) is strictly log-concave at every interior position.
Equivalently,
\begin{equation}\label{eq:logconcavity}
 b_j^2>b_{j-1}b_{j+1}\qquad(1\le j\le r-1).
\end{equation}
Consequently the coefficient sequence is unimodal. When \(r<2\), the
strict inequalities have no interior indices and are vacuous.
\end{corollary}
\begin{proof}
For real \(0\le y<r\), define
\[
 R_L(y)=\frac{(r-y)(L+y)}{(y+1)(L+r+y+1)}.
\]
Since \(L\ge1\),
\[
 \frac{d}{dy}\log R_L(y)
 =-\frac1{r-y}+\frac1{L+y}-\frac1{y+1}
  -\frac1{L+r+y+1}<0.
\]
Indeed, the middle two terms sum to a nonpositive number and the other
two are negative. The exceptional ratio is \(R_n(y)\) multiplied by
\((n+2y+2)/(n+2y)\), another positive strictly decreasing function.
Thus the positive adjacent ratios \(b_{j+1}/b_j\) strictly decrease
in both cases. This is precisely \eqref{eq:logconcavity}. Reversing the
sequence, as in \eqref{eq:rho-b}, preserves strict log-concavity.
A sequence with decreasing positive adjacent ratios increases up to at
most one transition and then decreases, giving unimodality.
\end{proof}

\paragraph{A useful boundary: not real-rootedness.}
The conclusion cannot be upgraded to real-rootedness of every residual
polynomial. Example~\ref{ex:k4n2} gives
\(\rho_4(2,t)=2+8t+10t^2\), whose discriminant is
\(8^2-4\cdot2\cdot10=-16\). The coefficient inequalities are therefore
not being deduced from real zeros; the adjacent-ratio argument is essential.

\paragraph{A terminating hypergeometric form.}
As an optional compact encoding, in the regular cases
\[
 \sum_{j=0}^r b_jz^j
 =b_0\,{}_2F_1(-r,L;L+r+1;-z).
\]
This is simply the ratio formula written as a finite sum, so no convergence
issue arises. In the exceptional case, take the same expression with
\(L=n\) and apply \(1+(2/n)z\,d/dz\). The resulting coefficient
multiplier \((n+2j)/n\) produces \eqref{eq:ratioExceptional}.
The product formula, rather than this notation, remains our primary
evaluation.

\section{Polynomial dependence and coefficient asymptotics}\label{sec:poly}
For the rest of the article, \(u=1\). The two-weight formula was useful
because it made the surviving terms easy to separate; no further
specialization is needed to extract the conjectured generating functions.

\begin{theorem}[Parity polynomials]\label{thm:parity}
Fix \(k\ge1\), a parity of \(n\), and an admissible deletion index.
On that parity class, \(C_{k,a}(n)\) is a polynomial in \(n\) of exact
degree \(d_k=\lfloor k^2/4\rfloor\). More precisely, with
\(a=\eps+2j\),
\begin{equation}\label{eq:leading}
 C_{k,\eps+2j}(n)
 =K_k\binom rj n^{d_k}+O_k(n^{d_k-1}),
\end{equation}
where the leading constant is the same on both parity classes, and
\begin{align}
 K_{2r}&=
 \frac{\displaystyle\left(\prod_{\ell=1}^{r-1}\ell!\right)^2}
 {\displaystyle 2^r\prod_{\ell=0}^{2r-1}\ell!},
 &&r\ge1,\label{eq:Keven}\\
 K_{2r+1}&=
 \frac{\displaystyle r!\left(\prod_{\ell=1}^{r-1}\ell!\right)^2}
 {\displaystyle 2^r\prod_{\ell=0}^{2r}\ell!},
 &&r\ge0.\label{eq:Kodd}
\end{align}
For \(r=0\), the numerator product in \eqref{eq:Kodd} is understood
as empty, giving \(K_1=1\).
\end{theorem}
\begin{proof}
On a fixed parity class, every node is \(n/2\) plus a constant.
The node differences in \eqref{eq:VE} and \eqref{eq:VO} are therefore
constant, the pairwise sum factors have slope 1, and each odd column
factor \(O_j+1\) has slope \(1/2\). Thus the degree is
\[
 \binom e2+\binom o2+o
 =\begin{cases}r^2,&k=2r,\\r(r+1),&k=2r+1,\end{cases}
\]
with positive leading coefficient.

The ordinary Vandermonde product for the integer nodes \(0,1,\ldots,r\)
with \(j\) deleted is
\[
 \frac{\prod_{\ell=1}^{r}\ell!}{j!(r-j)!}
 =\binom rj\prod_{\ell=1}^{r-1}\ell!.
\]
For even \(k\), the other parity block has \(r\) consecutive nodes,
so the product of all ordinary differences is
\(\binom rj(\prod_{\ell=1}^{r-1}\ell!)^2\).
For odd \(k\), the other block has \(r+1\) consecutive nodes,
contributing one extra factor \(r!\). In both cases the odd column
factors contribute \(2^{-r}\) to the leading coefficient. The
factorial denominators combine to
\(\prod_{\ell=0}^{k-1}\ell!\). This gives
\eqref{eq:Keven}--\eqref{eq:Kodd}, independently of the parity of \(n\).
\end{proof}

\begin{corollary}[Binomial coefficient profile]\label{cor:asymptotic}
For fixed \(k\ge1\),
\begin{equation}\label{eq:rhoasymp}
 \rho_k(n,t)=K_k n^{d_k}(1+t)^r+O_k(n^{d_k-1})
\end{equation}
coefficientwise as \(n\to\infty\). In particular, if
\(p_{k,n}(\ell)=[t^\ell]\rho_k(n,t)/\rho_k(n,1)\), then
\[
 p_{k,n}(\ell)\longrightarrow 2^{-r}\binom r\ell
 \qquad(0\le\ell\le r).
\]
For each fixed positive real \(t\),
\begin{equation}\label{eq:deltaasymp}
 \delta_k(n,t)=(-1)^{k\binom n2}t^{e_k(n)}
 \left(K_kn^{d_k}(1+t)^r+O_{k,t}(n^{d_k-1})\right).
\end{equation}
\end{corollary}
\begin{proof}
Sum \eqref{eq:leading} in \eqref{eq:rho-b} and use the binomial theorem.
The normalization at \(t=1\) has leading coefficient \(2^rK_k>0\).
The determinant asymptotic follows from \eqref{eq:c13}.
\end{proof}
The first leading constants are
\[
 K_1=1,\qquad K_2=\frac12,\qquad K_3=\frac14,
 \qquad K_4=\frac1{48},\qquad K_5=\frac1{576}.
\]
There is also an exact, rather than merely asymptotic, way to obtain
all lower-order corrections. Expand the finite products
\eqref{eq:VE}--\eqref{eq:VO} on each parity class. The expansion
terminates at degree \(d_k\); there is no unproved remainder estimate
or analytic interchange of infinite expansions. The two parity classes
can have different subleading terms even though their leading terms agree.

The boundary \(n=0\) is compatible with these same parity polynomials.
In the minor expansion, every shifted prefix of positive length contains
\(U_{-1}=0\), while the unshifted term has indices \(0,\ldots,k-1\)
and determinant 1. Equivalently, at \(n=0\),
\(C_{k,0}(0)=1\) and all other admissible coefficients are zero.
This observation is important: the generating function below includes
its correct constant term, without a separately added correction.

\section{Rational generating functions and recurrences}\label{sec:gf}
\begin{theorem}[The prescribed denominators]\label{thm:denominators}
The generating functions \(F_k(x,t)\) have the rational forms
\eqref{eq:c14even} and \eqref{eq:c14odd}, with numerators
\(A_k(x,t)\in\Z[t,x]\). At this stage the numerator degrees are
bounded by the denominator degrees minus 1; their exact smaller degrees
will follow in Section~\ref{sec:reciprocity}.
\end{theorem}
\begin{proof}
For \(k=2r\), put \(d=r^2\). The product formula gives
\[
 t^{-rn}\delta_{2r}(n,t)
 =\sum_{j=0}^r C_{2r,2j}(n)t^{-j}.
\]
By Theorem~\ref{thm:parity}, the right side is a degree-\(d\)
polynomial on each parity class of \(n\), with the same leading
coefficient on the two classes. It can therefore be written
\[
 P(n)+(-1)^nR(n),\qquad \deg P\le d,\quad\deg R\le d-1,
\]
where the polynomial coefficients lie in \(\Q(t)\). For a polynomial
\(p(n)\) of degree at most \(s\), the series
\(\sum_{n\ge0}p(n)y^n\) has denominator dividing
\((1-y)^{s+1}\). One proof is to expand \(p(n)\) in the basis
\(\binom nj\) and use
\(\sum_{n\ge0}\binom nj y^n=y^j/(1-y)^{j+1}\).
Thus the denominator of \(F_{2r}\) divides
\[
 (1-t^r x)^{d+1}(1+t^r x)^d
 =(1-t^r x)(1-t^{2r}x^2)^d,
\]
which is precisely \eqref{eq:c14even}.

For \(k=2r+1\), put \(d=r(r+1)\). The sign in
\eqref{eq:main} is \((-1)^m\) at both \(n=2m\) and \(n=2m+1\).
More explicitly,
\begin{align*}
 \delta_k(2m,t)&=(-t^k)^m
       \sum_{j=0}^r C_{k,2j}(2m)t^{-j},\\
 \delta_k(2m+1,t)&=(-t^k)^m
       \sum_{j=0}^r C_{k,1+2j}(2m+1)t^{r-j}.
\end{align*}
Each displayed sum is a polynomial in \(m\) of degree at most \(d\).
Generating the even and odd subsequences with \(y=-t^kx^2\) gives
a denominator dividing \((1+t^kx^2)^{d+1}\), proving
\eqref{eq:c14odd}.

The resulting rational functions are proper in \(x\), as is immediate
from the binomial-basis expressions. Finally, the prescribed denominator
has constant term 1 and lies in \(\Z[t,x]\). Multiplying it by the
original series in \(\Z[t][[x]]\) gives coefficients in \(\Z[t]\).
Rationality says that all but finitely many coefficients vanish, so
\(A_k\in\Z[t,x]\), not merely \(\Q(t)[x]\).
\end{proof}

\begin{corollary}[Constant-coefficient recurrences]\label{cor:recurrence}
Let \(\mathsf E f(n)=f(n+1)\). For all \(n\ge0\),
\begin{align}
 (\mathsf E-t^r)^{r^2+1}(\mathsf E+t^r)^{r^2}
       \delta_{2r}(n,t)&=0,\label{eq:receven}\\
 (\mathsf E^2+t^{2r+1})^{r^2+r+1}
       \delta_{2r+1}(n,t)&=0.\label{eq:recodd}
\end{align}
\end{corollary}
\begin{proof}
The factor \(\mathsf E-\lambda\) lowers the polynomial degree of a
sequence \(\lambda^np(n)\) by one. Apply this observation to the even
case in the preceding proof. In the odd case, apply the same argument
to each parity subsequence with \(\mathsf E^2\) in place of
\(\mathsf E\).
\end{proof}
These are valid annihilating recurrences, without a claim of minimality.
For example, \(\delta_1(n,t)+t\delta_1(n-2,t)=0\) for \(n\ge2\),
while the even-shift order is \(2r^2+1\), independent of the matrix
size \(n\).

\section{Negative indices, exact degrees, and reciprocity}\label{sec:reciprocity}
Rationality alone does not explain the numerator degrees in Conjecture 14,
or the reversal symmetries in Conjecture 15. Both come from a reflection
of the same deleted-index products. The negative indices used here are
an algebraic continuation of the coefficient sequence, not matrices of
negative size.

\subsection{A canonical bilateral continuation}
Extend \(U_d\) to every integer index by the recurrence
\eqref{eq:U}, working with invertible \(t\). It immediately yields
\begin{equation}\label{eq:Unegative}
 U_{-d-2}(z;t)=-t^{-d-1}U_d(z;t)\quad(d\ge0),
 \qquad U_{-1}=0.
\end{equation}
The coefficient identities \eqref{eq:evenentry}--\eqref{eq:oddentry}
continue to hold for every integer \(m\), if binomial coefficients
with negative upper argument are interpreted polynomially:
\(\binom zs=z(z-1)\cdots(z-s+1)/s!\).
One may check this directly from \eqref{eq:Unegative}; the two
quadratic factorizations \eqref{eq:quadE}--\eqref{eq:quadO} are
polynomial identities and remain valid as well.

Now use the right side of \eqref{eq:small}, with
\(P_d=U_d-U_{d-1}\), to define \(\delta_k(n,t)\) for all integer
\(n\). The prefix expansion, parity count, Vandermonde evaluation,
and sign calculation continue verbatim. Thus \eqref{eq:main}, with
\(u=1\) and the product \eqref{eq:C}, remains valid at all integer
\(n\), although its coefficients need not be positive when \(n<1\).
On each parity class the product is exactly the polynomial already
identified in Theorem~\ref{thm:parity}. Consequently this is the
exponential-polynomial continuation determined by the nonnegative values;
it is not a new choice of values made to force a symmetry.

\begin{lemma}[A negative-index gap]\label{lem:gap}
For every \(k\ge1\),
\begin{equation}\label{eq:gap}
 \delta_k(-1,t)=\delta_k(-2,t)=\cdots
 =\delta_k(-(k-1),t)=0.
\end{equation}
For \(k=1\), the displayed list is empty.
\end{lemma}
\begin{proof}
Let \(n=-h\), \(1\le h\le k-1\). Before a deletion, the degree
interval in the prefix expansion contains \(-2,-1,0\). If a degree
other than \(-1\) is deleted, the zero column \(U_{-1}\) remains.
If \(-1\) is deleted, the proportional columns
\(U_{-2}=-t^{-1}U_0\) and \(U_0\) both remain. Every term vanishes.
\end{proof}

\subsection{Reflection of the product coefficients}
\begin{lemma}[Coefficient reflection]\label{lem:Creflection}
For an admissible index and any integer \(m\),
\begin{equation}\label{eq:Creflection}
 C_{k,a}(-m)=(-1)^{d_k}C_{k,k-a}(m-k).
\end{equation}
\end{lemma}
\begin{proof}
The map \(d\mapsto-d-2\) sends the selected degree set
\(S_{k,-m,a}\) to \(S_{k,m-k,k-a}\), with the order reversed.
It sends an even node \(E\) to \(-E-1\), and an odd node \(O\)
to \(-O-2\). Each even pairwise sum factor changes sign, as does
each odd pairwise sum factor and each odd column factor.
Reversing the order makes the ordinary difference factors agree with
their counterparts. The total sign is therefore
\[
 (-1)^{\binom e2+\binom o2+o}=(-1)^{d_k}.
\]
The factorial denominators are unchanged, proving the identity.
\end{proof}

\begin{theorem}[Bilateral determinant reflection]\label{thm:reflection}
Put
\(s_k=(-1)^{\lceil k/2\rceil}\) and \(B_k=\binom{k+1}{2}\).
Then for every integer \(m\),
\begin{equation}\label{eq:reflection}
 \boxed{\displaystyle
 \delta_k(-m,t)=s_k t^{-B_k}\delta_k(m-k,t^{-1}).}
\end{equation}
In particular,
\begin{equation}\label{eq:firstnonzero}
 \delta_k(-k,t)=s_k t^{-B_k}\ne0.
\end{equation}
\end{theorem}
\begin{proof}
Apply Lemma~\ref{lem:Creflection} termwise in \eqref{eq:main} and
replace \(a\) by \(k-a\). This bijects the two admissible index sets.
The exponents of \(t\) match because
\[
 -B_k-\frac{k(m-k)-(k-a)}{2}=\frac{-km-a}{2}.
\]
The ratio of the two determinant signs has exponent, modulo 2,
\[
 d_k+\frac{k\{m(m+1)-(m-k)(m-k-1)\}}2
 \equiv d_k-\frac{k^2(k+1)}2
 \equiv\left\lceil\frac k2\right\rceil.
\]
For the last congruence, substitute \(k=2r\) or \(k=2r+1\).
This gives \eqref{eq:reflection}. Setting \(m=k\) uses
\(\delta_k(0,t^{-1})=1\) and gives \eqref{eq:firstnonzero}.
\end{proof}

\subsection{Reading a rational function at infinity}
We need an elementary link between this continuation and rational
functions. It is stated explicitly to justify the exact degree argument.

\begin{lemma}[Expansion at infinity]\label{lem:infinity}
Suppose \(f_n\) is a finite sum of terms \(p(n)\lambda^n\), with
\(p\) a polynomial and \(\lambda\ne0\), defined for every integer
\(n\) over a field of characteristic zero. Then the rational function
\(F(x)=\sum_{n\ge0}f_nx^n\) has Laurent expansion
\begin{equation}\label{eq:infinity}
 F(x)=-\sum_{m\ge1}f_{-m}x^{-m}
\end{equation}
at infinity.
\end{lemma}
\begin{proof}
For \(f_n=\lambda^n\), expand
\[
 \frac1{1-\lambda x}
 =-\frac{(\lambda x)^{-1}}{1-(\lambda x)^{-1}}
 =-\sum_{m\ge1}\lambda^{-m}x^{-m}.
\]
Apply \((x\,d/dx)^j\) to obtain the assertion for
\(f_n=n^j\lambda^n\), and use linearity.
\end{proof}
For our sequences the lemma applies over an algebraic extension of
\(\Q(t)\) if necessary. In the even case the bases are \(t^r\) and
\(-t^r\). In the odd case one can use the two square roots of
\(-t^{2r+1}\). All final identities lie in \(\Q(t)(x)\), so this
intermediate field extension introduces no restriction.

\begin{theorem}[Exact numerator degrees]\label{thm:degrees}
The numerators in Theorem~\ref{thm:denominators} have exactly the
degrees asserted in Conjecture 14:
\[
 \deg_xA_{2r}=2r(r-1)+1,\qquad
 \deg_xA_{2r+1}=2r^2+1.
\]
\end{theorem}
\begin{proof}
By Lemma~\ref{lem:infinity}, the negative-index gap and its first
nonzero successor give
\begin{equation}\label{eq:Fleadinginfty}
 F_k(x,t)=-s_k t^{-B_k}x^{-k}+O(x^{-k-1})
 \qquad(x\to\infty),
\end{equation}
as a formal Laurent expansion. Thus, writing \(F_k=A_k/Q_k\) with
the prescribed denominator, we have
\(\deg_xA_k=\deg_xQ_k-k\). The two denominator degrees are
\(2r^2+1\) and \(2(r^2+r+1)\), giving the required formulas.
No assertion that \(A_k\) and \(Q_k\) are relatively prime is needed:
the degree difference describes the rational function at infinity even
when both polynomials contain a common factor.
\end{proof}

\subsection{The reciprocal identities}
\begin{theorem}[Generating-function and numerator reciprocity]\label{thm:reciprocity}
As rational functions,
\begin{equation}\label{eq:Freciprocity}
 F_k(x,t)=-s_k t^{-B_k}x^{-k}F_k(x^{-1},t^{-1}).
\end{equation}
Consequently \(A_k\) satisfies \eqref{eq:c15even} and
\eqref{eq:c15odd}, proving Conjecture 15.
\end{theorem}
\begin{proof}
Insert \eqref{eq:reflection} in \eqref{eq:infinity}:
\[
 F_k(x,t)=-s_k t^{-B_k}\sum_{m\ge1}
                 \delta_k(m-k,t^{-1})x^{-m}.
\]
For \(1\le m<k\), every summand vanishes by the negative-index gap.
Reindex the remaining terms by \(n=m-k\) to obtain
\eqref{eq:Freciprocity}. Equality of the Laurent expansions implies
equality of the rational functions.

For even \(k=2r\), direct factorization of the denominator gives
\begin{equation}\label{eq:Qreceven}
 Q_{2r}(x^{-1},t^{-1})
 =(-1)^{r^2+1}t^{-r-2r^3}x^{-1-2r^2}Q_{2r}(x,t).
\end{equation}
Combine this with \eqref{eq:Freciprocity} and
\(B_{2r}=2r^2+r\). The signs cancel, and
\[
 A_{2r}(x^{-1},t^{-1})
 =t^{-2r^2(r-1)}x^{-2r(r-1)-1}A_{2r}(x,t).
\]
This is \eqref{eq:c15even}.
For odd \(k=2r+1\), put \(d=r(r+1)\). Now
\begin{equation}\label{eq:Qrecodd}
 Q_k(x^{-1},t^{-1})=t^{-k(d+1)}x^{-2(d+1)}Q_k(x,t).
\end{equation}
Since \(B_k=k(r+1)\), substitution gives
\[
 A_k(x^{-1},t^{-1})=(-1)^r
          t^{-kr^2}x^{-2r^2-1}A_k(x,t),
\]
which is \eqref{eq:c15odd}.
\end{proof}

\begin{corollary}[Leading numerator coefficients]\label{cor:Aleads}
All numerators have constant term \(A_k(0,t)=1\). Moreover,
\begin{align*}
 [x^{2r(r-1)+1}]A_{2r}(x,t)&=t^{2r^2(r-1)},\\
 [x^{2r^2+1}]A_{2r+1}(x,t)&=(-1)^r t^{(2r+1)r^2}.
\end{align*}
Their degrees in \(t\) are respectively \(2r^2(r-1)\) and
\((2r+1)r^2\).
\end{corollary}
\begin{proof}
The constant term comes from \(Q_k(0,t)=F_k(0,t)=1\).
Compare the opposite extreme coefficients in the reciprocity identity
to obtain the leading terms. Since \(A_k\) is a polynomial in \(t\),
reciprocity also says that no \(t\)-exponent can exceed the stated
reciprocity exponent. The leading terms attain that bound.
\end{proof}
All uses of \(t^{-1}\) above take place over \(\Q(t)\).
The resulting numerator identities are polynomial identities, so they
extend to \(t=0\) after their displayed monomial factors are multiplied
out. The rational description of \(F_k\), whose denominator has
constant term 1 in \(x\), is likewise valid as a formal generating
function under every parameter specialization.

\section{Examples and computational consequences}\label{sec:examples}
\subsection{Small numerator polynomials}
Finite convolution of the product formula with the denominator produces
\begin{align*}
 A_1(x,t)&=1+x,\\
 A_2(x,t)&=1+x,\\
 A_3(x,t)&=1+(1+2t)x-t^2(2+t)x^2-t^3x^3,\\
 A_4(x,t)&=1+(1+3t+t^2)x+t^2(1+t)(1+4t)x^2\\
 &\quad+t^4(1+t)(4+t)x^3
       +t^6(1+3t+t^2)x^4+t^8x^5.
\end{align*}
These agree with the small examples in the original paper
\cite[Section 5]{Cigler2021}. The signs of the numerator coefficients
are not uniformly positive: \(A_3\) already has negative coefficients.
The positivity theorem concerns \(\rho_k\), not \(A_k\).
The first few degree data are
\[
\begin{array}{c|rrrrrrrr}
 k&1&2&3&4&5&6&7&8\\ \hline
 d_k&0&1&2&4&6&9&12&16\\
 \deg_x Q_k&2&3&6&9&14&19&26&33\\
 \deg_x A_k&1&1&3&5&9&13&19&25.
\end{array}
\]
The accompanying \texttt{data/numerators.json} contains all numerator
polynomials through \(k=10\), with integer coefficient triples
\((i,j,c)\) representing \(cx^it^j\). It also contains the prescribed
denominators in the same format, making the data independent of a
particular computer-algebra syntax.

\subsection{Residual polynomials}
For example, the product formula gives
\begin{align*}
 \rho_2(5,t)&=3+3t,\\
 \rho_3(6,t)&=12+16t,\\
 \rho_4(7,t)&=120+270t+150t^2,\\
 \rho_5(8,t)&=1650+4400t+3025t^2,\\
 \rho_6(9,t)&=55055+195195t+231231t^2+91091t^3.
\end{align*}
The coefficients can be large even when the polynomial degree is small.
This is precisely the situation in which an explicit product is more
useful than expansion of a large determinant.

\subsection{An exact coefficient algorithm}
The evaluation has two natural implementations. The most transparent
one independently evaluates \eqref{eq:C} for each of the \(r+1\)
terms. A faster one evaluates \(b_0\) once and then uses the rational
adjacent ratios of Theorem~\ref{thm:ratios}. Every resulting update is
an integer, so exact division is an internal consistency check.

In arithmetic-operation counts, the first coefficient uses \(O(k^2)\)
pairwise factors and the remaining coefficients use \(O(k)\) updates.
The determinant dimension \(n\) no longer determines a matrix-elimination
cost. This is an arithmetic count, not a unit-cost assertion for
arbitrarily large integers: coefficient bit lengths grow with \(n\)
and \(k\), and a numerical value of \(t^{e_k(n)}\) may itself require
many bits. The reference implementation provides both coefficient
algorithms and checks them against each other.

For generating functions, write \(Q_k(x,t)=\sum_{j=0}^{D}q_j(t)x^j\)
and let \(N=D-k\). Then
\begin{equation}\label{eq:Analgorithm}
 A_k(x,t)=\sum_{n=0}^{N}
 \left(\sum_{j=0}^{\min(n,D)}q_j(t)\delta_k(n-j,t)\right)x^n.
\end{equation}
Thus \(N+1\) explicitly known determinant values suffice to construct
the entire numerator. The proof establishes that every later coefficient
of \(Q_kF_k\) vanishes; a finite computation is not being used to guess
that vanishing.

\paragraph{Why this particular attack succeeds.}
The decisive structural feature is not a large computation. It is the
short connection \(P_m=U_m-uU_{m-1}\). Selecting terms in adjacent
columns produces duplicate columns except for a single switch point.
The problem is reduced to choosing which one of \(k+1\) consecutive
indices is omitted. Parity then cuts the count roughly in half, and the
Chebyshev coefficient arrays have quadratic Vandermonde evaluations.
Positivity, parity-polynomial growth, and reflection are three aspects
of that same product, rather than unrelated lucky patterns.

\section{Exact verification and reproducibility}\label{sec:verification}
The proof above is the basis for the general statements. The accompanying
computations serve a different purpose: they check indexing, signs,
normalization, boundary values, examples, and the implementation through
independent routes. No floating-point comparison is used in the tests.

The standard-library suite completed \textbf{32,666 exact checks}.
Its principal groups are summarized in Table~\ref{tab:checks}; the full
machine-readable record is \texttt{data/verification.json}.
\begin{table}[htbp]
\centering
\small
\begin{tabular}{@{}p{0.63\textwidth}r@{}}
\toprule
Check family & Number\\
\midrule
Independent shifted Hankel determinants: grid & 1,092\\
Independent shifted Hankel determinants: seeded cases & 200\\
Independent unshifted Hankel determinants & 91\\
Positive individual coefficients & 10,080\\
Adjacent-ratio identities & 8,640\\
Strict log-concavity inequalities & 7,260\\
Agreement of direct and ratio-based coefficient algorithms & 1,440\\
Parity-polynomial degree and leading-term checks & 70\\
Negative-index reflection identities & 2,352\\
Negative-index gap values & 360\\
Vanishing generating-function tail coefficients & 984\\
Numerator reciprocity and exact degree, rational parameters & 96\\
Non-real-rooted counterexample & 1\\
\midrule
Total & 32,666\\
\bottomrule
\end{tabular}
\caption{Exact checks performed by \texttt{code/verify.py}. The count
records individual assertions, not independent theorems.}\label{tab:checks}
\end{table}

The determinant grid uses \(1\le k\le12\), \(0\le n\le12\), and
seven integer pairs \((u,t)\), including zero and negative weights.
The seeded additional cases have \(1\le k\le20\) and
\(0\le n\le18\). Each direct determinant is evaluated from the
ballot-sum moments by fraction-free Bareiss elimination with pivoting;
it does not use the small orthogonal-polynomial determinant or the
product formula internally. At \(t=0\), the comparison is made with
\(t^{\binom n2}\delta_k(n;u,t)\), never by dividing a numerical
determinant by zero.

The coefficient inequalities and ratios cover
\(1\le k\le24\), \(1\le n\le60\).
Finite differences check both parity-polynomial degrees and their exact
leading terms through \(k=10\).
The bilateral and generating-function checks use \(1\le k\le16\)
and \(t\in\{1,2,-1/2\}\). Reflection is checked for
\(0\le m\le48\); generating-function coefficients are checked through
\(x^{\deg Q_k+12}\).

A second suite adds \textbf{24 independent bivariate symbolic determinant
identities} for \(1\le k\le6\), \(1\le n\le4\), retaining both
\(u\) and \(t\) as indeterminates. It also checks 135 symbolic tail
coefficients and the complete symbolic reciprocity identity for each
\(1\le k\le10\). The symbolic determinant calculations use SymPy;
the latter generating-function calculations use sparse integer polynomial
arithmetic and retain symbolic \(t\). The saved report identifies the
actual runtime and SymPy versions.

\subsection{Running and rebuilding}
From the extracted archive directory, run
\begin{lstlisting}
python code/verify.py
python -m pip install -r requirements.txt
python code/verify_symbolic.py
pdflatex -interaction=nonstopmode -halt-on-error article.tex
pdflatex -interaction=nonstopmode -halt-on-error article.tex
\end{lstlisting}
The first test needs only Python 3.9 or later. The second test requires
SymPy; the tested version is recorded in \texttt{requirements.txt}.
The \LaTeX{} source uses ordinary packages distributed with TeX Live or
MiKTeX. All references are included in the source, so there is no BibTeX
or external bibliography-generation step. Shell and PowerShell build
scripts are included as conveniences.

The principal artifacts are \texttt{article.tex}, \texttt{article.pdf},
\texttt{code/hankel.py}, the two verification scripts, and the data files.
\texttt{data/rho\_coefficients.csv} supplies coefficients for
\(1\le k\le12\), \(1\le n\le20\). No downloaded paper or font file
is needed to build or test the package.

\section{Conclusion and scope}\label{sec:conclusion}
The explicit deleted-index product resolves the three stated mathematical
claims as follows. Positivity and the precise support prove Conjecture 13.
Parity-polynomial dependence gives the predicted denominators, while the
negative-index gap supplies the exact numerator degrees, completing
Conjecture 14. Reflection of the even and odd nodes yields the rational
functional equation and hence Conjecture 15. The adjacent-ratio formulas
further prove strict log-concavity, and the leading products give a
binomial coefficient profile with explicit constants.

The scope is exactly these three source conjectures and the additional
results stated here. No conclusion is claimed about the other conjectures
in the source paper, universal real-rootedness, or a reduced denominator
after every parameter specialization. Historical priority remains subject
to the qualification in Section~\ref{sec:problem}; the proof and the
verification artifacts are provided for independent examination.

\begin{thebibliography}{9}
\small
\setlength{\itemsep}{3pt}
\setlength{\parskip}{0pt}
\bibitem{Cigler2021}
Johann Cigler.
\emph{Hankel determinants of middle binomial coefficients and conjectures
for some polynomial extensions and modifications}.
arXiv:2111.14492v3, revised 30 December 2021.
In particular, Section 5, pp. 19--21, Conjectures 13--15.
\url{https://arxiv.org/abs/2111.14492v3}.

\bibitem{CK2020}
Johann Cigler and Christian Krattenthaler.
\emph{Hankel determinants of linear combinations of moments of orthogonal
polynomials}.
arXiv:2003.01676, 2020.
\url{https://arxiv.org/abs/2003.01676}.

\bibitem{K2021}
Christian Krattenthaler.
\emph{Hankel determinants of linear combinations of moments of orthogonal
polynomials, II}.
arXiv:2101.04225, 2021.
\url{https://arxiv.org/abs/2101.04225}.

\bibitem{OEIS1405}
OEIS Foundation Inc.
\emph{The On-Line Encyclopedia of Integer Sequences}, entry A001405:
\(a(n)=\binom n{\lfloor n/2\rfloor}\).
Accessed 19 September 2026.
\url{https://oeis.org/A001405}.
\end{thebibliography}
\end{document}
